\documentclass{article}
\usepackage[utf8]{inputenc}
\usepackage[spanish]{babel}

\title{Bitácora}    
\begin{document}

\maketitle

\section{17 de agosto de 2026, clase de tópicos}



**Preguntas del día de hoy**
*¿Qué entiendo actualmente por investigación científica en física?*
-Hoy entiendo actualmente a una investigación científica en física como una hipótesis propuesta sobre algo físico que ocurre en la naturaleza, la cuál debe ser modelada matemáticamente  y  verificada experimentalmente, así comprobandola o no. 

*¿Qué diferencias identifico entre resolver un ejercicio de física  desarrollar una investigación?*
Resolver un ejercicio de física es resolver un problema creado anteriormente con el propósito de mejorar la comprensión del asusnto de quien lo hace, mientras que desarrollar una investigación científica es buscar comprender algo desconocido,  o mejorar la comprención de un fenómeno físico.

*¿Qué espero observar o aprender durante el IX workshop de Ciencias Físicas?*
Espero ver que se realiza en áreas que me llaman la atención de la física, para hacerme una mejor idea del trabajo real y de como es la investigación en dicha área.

-Hoy se habló de lo que se va a realizar en el workshop, junto con las preguntas que se tienen que hacer en este para la bitácora, la cual luego se tendrá que escribir en $LaTex$ y alojar en github.

Tarea!!! Hacer presentación con netebooklm, con tres aplicaciones.(intentar hacer sola).


\section{18/08/26 workshop}

Preguntas
1-Que problema o pregunta de investigación aborda?
2-por qué resulta interesante o relevante ?
3-que método utiliza: teoría, simulación , análisis de datos u otros?
4-cuál es el resultado principal?
5-que figura o resultado considera principalmente importante y que muestra?
6- que aspectos del trabajo no comprendió completamente?
7-considera que algún resultado del artículo podría ser reproducido con los conocimientos y herramientas disponibles actualmente en la carrera? Por qué?



RELATIVIDAD GENERAL Y GEOMETRÍA DIFERENCIAL

1-que busca explicar la relatividad general?
2-por que la gravedad es  una de las fuerzas fundamentales
3-teoría
4-El cómo se describe matemáticamente la relatividad general
5-la ecuación de la acción einstein-hilbert , una de calcular la acción de la relatividad general.
6-No comprendí completamente conceptos matemáticos como por ejemplo los tensores.
7-creo que actualmente me faltan bases matemáticas para comprender la relatividad general.


Ecuaciones BPS en el modelo Sigma no lineal y condensado de Bose-Einstein

1-El problema de las ecuaciones no lineales, por su dificultad, de los solitones debido a su topología.
2- Resulta interesante debido a la única tipología de los solitones
3-utiliza la teoría
4- El formalismo BPS permite resolver ecuaciones de manera lineal.
5- considere importante la figura de las ecuaciones BPS para GPE, ya que es un resultado nuevo que muestra las aplicaciones de las ecuaciones BPS.
6- No comprendi las matemáticas detrás de esta presentación, tampoco la física de la materia.
7- Consideró que esta presentación es muy avanzada a los conocimientos que tengo en este momento, por lo que no creo que sea capaz de reproducir nada de esto.


\section{19/08/26 workshop}

Preguntas
1-Que problema o pregunta de investigación aborda?
2-por qué resulta interesante o relevante ?
3-que método utiliza: teoría, simulación , análisis de datos u otros?
4-cuál es el resultado principal?
5-que figura o resultado considera principalmente importante y que muestra?
6- que aspectos del trabajo no comprendió completamente?
7-considera que algún resultado del artículo podría ser reproducido con los conocimientos y herramientas disponibles actualmente en la carrera? Por qué?
8-preguntas que haría 









Hamiltoniano de campo medio

1-que ocurre cuando los sistemas están formados con partículas que interactúan a largas distancias?
2-porque no es despreciable la contribución de las partículas lejanas
3-teoría
4-El cómo se puede ocupar el modelo hamiltoneano para una guía de onda con átomos
5-el modelo de la situación física  hablado en 4
6-No comprendí completamente el como habla de la descripción de el movimiento de los átomos y varios parámetros.
7-consideró que aún no puedo responder a un resultado de esto con mis conocimientos actuales.


\section{23 de agosto 2026}

**23 de agosto 2026**

Me di cuenta que respondí las preguntas de la presentación en el workshop :(

Hoy haré las presentaciones de los papers que elegí. Estos son :
-*Noise Doppler-Shift Measurement of Airplane Speed  
-Searching nontrivial magnetic equilibria using the deflated Newton method  
-Magnetic field ‘flyby’ measurement using a smartphone’s magnetometer and accelerometer simultaneously

Preguntas a responder:
¿qué problema o pregunta de investigación aborda?;
¿por qué resulta interesante o relevante?;
¿qué método utiliza: experimento, teoría, simulación, análisis de datos u otro?;
¿cuál es el resultado principal?;
¿qué figura o resultado considera especialmente importante y qué muestra?;
¿qué aspectos del trabajo no comprendió completamente?;
¿considera que algún resultado del artículo podría ser reproducido con los conocimientos y
herramientas disponibles actualmente en la carrera?, ¿por qué?



-**Magnetic field ‘flyby’ measurement using a smartphone’s magnetometer and accelerometer simultaneously**

¿qué problema o pregunta de investigación aborda?;
Aborda el cómo obtener experimentalmente el campo magnético generadopor una bobina con el magnetrómmetro y el acelerómetro de un celular.

¿por qué resulta interesante o relevante?;
Porque muestra una interesante aplicación de como los sensores de los celulares permiten obtener medidas simultaneas de parámetros fisícos, lo que permite obtener experimentalmente el campo eléctrico, cosa que es normalmente medida con un sensor Hall.

¿qué método utiliza: experimento, teoría, simulación, análisis de datos u otro?;
Utiliza la experimentación principalmente, obteniendo el valor del campo magnético con los sensores del celular y una aplicación de este mismo, para posteriormente analizar los datos constrarstando los resultados con el modelo teórico basado en  la ley de Bio -Savart.

¿cuál es el resultado principal?
El resultado principal fue que el valor de la permeabilidad se encuentra dentro de los valores aceptables,  que se calculo con los resultados obtenidos experimentalmente, por lo que se concluyó que fue una manera bastante precisa de medir el campo magnético.

¿qué figura o resultado considera especialmente importante y qué muestra?;
Considero a la figura 2 especialmente importante, ya que muestra la gran concordancia entre los resultados experimentales y el modelo teórico, también se puede ver como la concordancia entre ambas es peor mientras la distancia es mas grande. 

¿qué aspectos del trabajo no comprendió completamente?;
No comprendo muy bien la ecuación de Bio-Stewart y los conceptos de electromagnetismo en general, aún no obtengo esos conocimientos.

¿considera que algún resultado del artículo podría ser reproducido con los conocimientos y
herramientas disponibles actualmente en la carrera?, ¿por qué?
Considero que este experimento podría ser reproducido con mis conocimientos actuales, ya que aun que no manejo al 100% los conceptos el experimento y el análisis posterior es suficientemente simple como para realizarlos con una investigación previa de los conceptos con un poco más de profundidad y realizarlo sin problemas. Mas opuede ser que las os materiales sean un poco más complicados de conseguir.




**-Noise Doppler-Shift Measurement of Airplane Speed**

¿qué problema o pregunta de investigación aborda?;
Aborda el cómo se puede medir la velocidad de un avión utilizando el efecto Doppler

¿por qué resulta interesante o relevante?;
Resulta relevante porque presentan una nueva, moderna y barata forma de hacer experimentos que involucran el efecto Doppler. 

¿qué método utiliza: experimento, teoría, simulación, análisis de datos u otro?;
Este artículo utiliza la experimentación, obteniendo el sonido de los aviones grabando su sonido con un micrófono o grabación, y posteriormente el análisis de datos utilizando un programa que analiza el sonido y realiza transformaciones de fourier obteniendo así su frecuencia, y mostrandose el efecto Doppler.

¿cuál es el resultado principal?;
El resultado principal es el cómo se puede comprobar experimentalmente el efecto Doppler que se produce con el sonido de los aviones, así mostrando una aplicación de este mismo, calculando la velocidad con las frecuencias obtenidas.

¿qué figura o resultado considera especialmente importante y qué muestra?

Considero a las figuras 2 y 3 especialmente inportantes, ya que la figura 2 muestra el espectrograma de la velocidad vs frecuencia después de la FFT  del sonido grabado, y la figura 3 muestra un zoom de la figura 2, en donde se ve el efecto Doppler, ya que emieza en 2,2 kHz y termina en 1,5 kHz.

¿qué aspectos del trabajo no comprendió completamente?
No comprendo bien las matemáticas de las FFT.


¿considera que algún resultado del artículo podría ser reproducido con los conocimientos y
herramientas disponibles actualmente en la carrera?, ¿por qué?
Considero que este artículo es completamente reproducible con los conocimientos disponibles actualemte, ya que utiliza ecuaciones ya vistas en física ondulatoria y programas y materiales facilmente obtenibles.

Uso de IA (LUMO)(avión)
prompt:
como cito una imagen en una presentacion en apa7
Type Journal Article Noise Doppler-Shift Measurement of Airplane Speed Costa Ivan F. Mocellin Alexandra The Physics Teacher 2007-09-01 45 6 356-358 10.1119/1.2768692 https://pubs.aip.org/pte/article/45/6/356/275475/Noise-Doppler-Shift-Measurement-of-Airplane-Speed 8/10/2026, 8:48:17 PM 0031-921X, 1943-4928 en DOI.org (Crossref) 8/10/2026, 8:48:17 PM 8/10/2026, 8:48:17 PM vagoberto This paper illustrates a new and practical experimental technique for studying the Doppler effect1–3 where the pitch variation4 of noise from a passing aircraft is used to calculate its speed.



-**Searching nontrivial magnetic equilibria using the deflated Newton method 

¿qué problema o pregunta de investigación aborda?;
Aborda el encontrar estados de equilibrio en sistemas  no lineares en donde hay presencia de simetría utilizando el método de newton deflactado. En este caso se estudia un sistema de rotores magnéticos macroscópicos.

¿por qué resulta interesante o relevante?;
Porque las herramientas presentadas en este artículo son fácilmente aplicables a otros sistemas dinámicos no lineales con simetrías discretas, como por ejemplo las redes de osciladores idénticos, utilizadas en la neurociencia matemática y otros contextos.

¿qué método utiliza: experimento, teoría, simulación, análisis de datos u otro?;
Utiliza un método experimental al principio para ver diversas texturas magnéticas. Lo principal del artículo es la simulación sistemas para buscar equilibrios estables con el método de Newton Deflactado. también contiene análisis de datos en donde traducen los resultados en códigos de colores.

¿cuál es el resultado principal?;
El resultado principal es que el método de newton deflactado se comportaba peor que el metodo de newton "ingenuo" para un numero moderado de dimensiones . Contrario a las expectativas, el metodo de la deflación no  ofrecia certeza de encontrar un gran numero de soluciones. También se tiene que el metodo de newton deflacionado puede enfrentar limitaciones en computadores con baja potencia

¿qué figura o resultado considera especialmente importante y qué muestra?
La figura más representativa sería la figura 8 , ya que muestra una comparación del método de newton normal y el método de newton deflactado. La figura a muestra el método de newton normal y el núcleo blanco serían las soluciones triviales. en la parte b se utiliza el método de newton deflactado y se ve el nucleo blanco desaparece y las soluciones no triviales aparecen de una forma mas compleja.En la parte c se deflactan solo los extados mas visitados por el método clásico.

¿qué aspectos del trabajo no comprendió completamente?
No comprendo el método de newton, ni el método de newton deflacionado, tampoco teoría de grupos , ni rotores magnéticos, realmente casi nada de este artículo.

¿considera que algún resultado del artículo podría ser reproducido con los conocimientos y
herramientas disponibles actualmente en la carrera?, ¿por qué?
Considero que ningún resultado obtenidop en este artículo es reproducible con mis conociminetos actuales.

USO de IA (GEMINI NOTEBOOK)
prompt:¿qué problema o pregunta de investigación aborda?(Se hizo casi omiso a casi todo lo que decia para poner solamente lo que consideraba proncipal); ¿qué método utiliza: experimento, teoría, simulación, análisis de datos u otro?;
que representan las figuras? ();
cual es la figura mas representativa de los resultados y por que
Se constrarrestó con el  anterior análisis del artículo.



\section{24 agosto 2026, clase de tópicos}

Presentamos las presentaciones

Preguntas a responder
1. ¿Qué problema o pregunta estudian?
2. ¿Qué método utilizan: experimento, teoría, simulación, datos?
3. ¿Cuál parece ser el resultado principal?
4. ¿Qué concepto o término no entendí?
5. ¿Qué pregunta le haría al expositor?

Presentación de Bruno y Jose owo
Una forma de enseñar el efecto doppler
1-Como medir el efecto Doppler con celulares
2-Experimentación y análisis de datos
3-La demostración del gráfico del efecto Doppler
4- todo muy entendible

-me faltó colocar los autores :(

Francisca y Benjamin
Video analysis to examine Kepler's laws of planetary motion
1-Analizan la segunda ley de kepler con instrumentos muy fáciles de conseguir
2-Experimentacion y análisis de datos


Analisis del efecto doppler en un expectrograma de un vehiculo avanzando
1-Como analizan el efecto Doppler de un vehiculo en la carretera
2-Experimentacion y análisis de datos
3-
4-El gráfico del expectrograma

Calcular cte g de gravitación con el efecto doppler
1- el titulo
2-Experimentación  y análisis de datos
3-
4-

Gabriel
Air resistance
1-Ve como el cae un objeto en caida parabolica
2-Experimentación 
3-
4-
5-

campo magnético
1- como obtener el campo magnetico con el magnetómetro y un acelerómetro
2- Experimentación y análisis de datos
3-La obtención del campo magnético
4-
5-

Cony
Aplicaciones del espectro sonoro
1-Como analizar los sonidos tubo
2-Experimentación y análisis de datos
3- Que el experimento sirve bien para las frecuencias altas
4-awa


Seba
Medicion de la velocidad del sonido mediante el analisis de los modos normales con un celular
1-Medir la velocidad del sonido con un microfono del celular
2-Experimentación y análisis de datos
3-Se obtuvo la velocidad del sonido con una buena presición
4-
5-

Fran
Simplificación de sistemas acoplados con la transformada de fourier
1- El título3
2-Experimentación y análisis de datos
3-La FFT simplifica mucho todo
4-
5-


\section{31/08/26}
Presentaciones

Preguntas a responder
1. ¿Qué problema o pregunta estudian?
2. ¿Qué método utilizan: experimento, teoría, simulación, datos?
3. ¿Cuál parece ser el resultado principal?
4. ¿Qué concepto o término no entendí?
5. ¿Qué pregunta le haría al expositor?

Juli
kepler
1-Simular y comprobar empiricamente la segunda ley de kepler utizando análisis de video.
2-Experimentación
3-Se comprobó la ey de kepler con un error debido a la friccion
4-
5-

Lahsen
Oscilaciones acopladas
1- Oscilaciones acopladas
2-Experimento
5-Como se verian los resultados


Lucas
1-Como calcular g con el efecto doppler
2-Experimentación, teoría y análisis de datos
3-Se obtuvo un valor muy parecido al valor real de g

Leo
1-La visualiazción directa de batidos mecánicos mediante un sistema masa-resorte
2- Experimentación
3-Se puede observar bien el fenómeno

Iñaki
1-Oscilaciones acopladas con celus (ver si los celulares son precisos)
2-Experimentación
3-El telefono es muy preciso

Rocío
1- Enseñar caida proyectil con celular
2-Experimentación y análisis de datos
3- 


Empezamos a organizar el trabajo, el nuestro será del avión.


\section{7/09/26}

Clase Tópicos

Términamos el formulario de planificación.
feedback :
-medir 3 ptos distintos
-probar mediciones entre cada avión
-entender la fft


\section{10/09/26}
Fuimos a buscar los aviones al humedal después del certamen. Habian muchos zancudos.


\section{11/09/26}

-Haciedo el código para el avión




\section{12/09/26}
Segui avanzando con el espectrograma
-buscar referencias de: -interpolacion cubica
-argumentar decibeles espectrograma
dB SPL


\section{16/09/26}

https://www.geeksforgeeks.org/machine-learning/cubic-spline-interpolation/
https://pythonnumericalmethods.studentorg.berkeley.edu/notebooks/chapter17.03-Cubic-Spline-Interpolation.html
https://docs.scipy.org/doc/scipy/reference/generated/scipy.interpolate.CubicSpline.html


\section{21/09/26}

Clases, vimos como hacer artículos científicos
Título "Reproducción del artículo blablalba..."
Afilicación: Departamento de Física, CFM, UDeC


7 de octubre workshop  física\\

Introducción\\
-que es el efecto doppler\\
aplicaciones\\
---> lo vamos a utilizar para esto\\

Métodos\\
que metodología se hizo\\

Resultados\\
Figuras  
----->Discusión
ir describirndo los resultados y discutirlos

Conclusiones
Resumen extendido 






\end{document}